% !TeX root = ../main.tex

\ustcsetup{
  keywords  = {交错磁性；自旋群对称性；非线性反常霍尔效应；贝利曲率多极子；CrSb},
  keywords* = {Altermagnetism; Spin-group symmetry; Nonlinear anomalous Hall effect; Berry curvature multipoles; CrSb}
}


\begin{abstract}
本文围绕交错磁体的前沿研究展开。系统阐释了交错磁性的物理内涵，并基于自旋群对称性原理，
详细论述了d波、g波及i波交错磁体的分类准则。在此基础上，本文探讨了已知交错磁体引入自旋轨道耦合（SOC）
后的新奇输运性质，
重点讨论了贝利曲率四极矩诱导的三阶非线性反常霍尔效应，并通过对称性分析揭示了晶格与自旋对称性对高阶非线性输运的内在约束机制。
特别地，基于CrSb单晶的旋转角度依赖测量和多波混频实验，验证了三阶横向非线性霍尔响应的各向异性及其瞬时宽带（${>}100\,\mathrm{kHz}$）特性。
\end{abstract}

\begin{abstract*}
This work focuses on recent advances in the study of altermagnets. First, the physical definition of altermagnetism is systematically elucidated, and the classification criteria for d-wave, g-wave, and i-wave altermagnets are detailed based on spin-group symmetry. Subsequently, the transport properties of known altermagnets are discussed, with particular emphasis on the nonlinear anomalous Hall effect induced by the interplay of quantum geometry and Berry curvature multipoles. Through symmetry analysis, the intrinsic constraints imposed by lattice and spin symmetries on high-order nonlinear transport are revealed. Notably, this paper discusses the experimental verification via multi-wave mixing measurements of the third-order nonlinear Hall effect in CrSb.
\end{abstract*}
